Esta sección presenta las expresiones aplicadas en el orden del flujo de trabajo. Se distingue entre procedimientos \textbf{ejecutados}, cuya evidencia está en scripts y salidas, y procedimientos \textbf{pendientes}, especificados para etapas futuras pero aún sin resultados para la matriz acordada. Los diagnósticos de una matriz anterior no se trasladan a la matriz candidata actual.

\subsection{Etapa 1. Verificación de estructura, llaves y cobertura}

Cada llave municipal se normaliza como texto de cinco dígitos, $c_i\in\{0,\ldots,9\}^5$. La consistencia departamental se evalúa con
\[
I_{D,i}=\mathbf{1}\{\texttt{codigo\_departamento}_i=\operatorname{substr}(c_i,1,2)\}.
\]
El control exige cero códigos duplicados y cero filas con $I_{D,i}=0$. La cobertura entre dos fuentes $A$ y $B$ se compara mediante la diferencia simétrica de conjuntos, $A\triangle B$. Para tablas con más de una fila municipal, la llave compuesta debe ser única: municipio--área--sexo, municipio--área--tipo de vivienda o municipio--área--grupo de edad, según el bloque.

La matriz de presencia define $P_{ij}=1$ cuando el código $i$ aparece en la fuente $j$, y $0$ en caso contrario. Las combinaciones de área, edad o tipo no publicadas por el diseño se cuentan como \emph{ausencias estructurales}, no como celdas faltantes ni como ceros imputados. El cruce histórico con DIVIPOLA actual identificaba 18 ANM. La verificación de la geografía de 2018 amplía la clasificación a veinte al incorporar Barrancominas y Mapiripana; se conservan separadas en la actualización descriptiva.

En reconciliaciones de conteos se calcula el residuo $e_i=T_i-\sum_{k=1}^{K}C_{ik}$, donde $T_i$ es el total publicado y $C_{ik}$ las categorías exhaustivas. Una fila viola la identidad si $|e_i|>0,5$ personas o viviendas. El rango de proporciones 0--1 se controla con $\mathbb{1}(x_i<0\ \text{o}\ x_i>1)$; para indicadores IPM se usa el rango [0,100]. La igualdad entre las columnas IPM Municipal e IPM se verifica con tolerancia $10^{-10}$. Los detalles por chequeo, hoja y tolerancia están en los CSV de verificación.

\subsection{Etapa 2. Construcción de indicadores candidatos}

Las expresiones usan conteos del CNPV 2018 del área \enquote{Total}. Para todo cociente se aplica $\operatorname{ratio}(a,b)=a/b$ si $b>0$; un denominador cero o indefinido genera NA y no un cero artificial.

\begin{singlespace}\small
\begin{longtable}{@{}p{0.15\textwidth}p{0.52\textwidth}p{0.23\textwidth}@{}}
\caption{Fórmulas operacionales de las candidatas}\label{tab:formulas_algebraicas}\\
\toprule
\textbf{Indicador} & \textbf{Fórmula operacional} & \textbf{Estado e interpretación}\\
\midrule
\endfirsthead
\toprule
\textbf{Indicador} & \textbf{Fórmula operacional} & \textbf{Estado e interpretación}\\
\midrule
\endhead
1HM & $H^{<15}_i/H_i$: hogares que reportan menores de 15 / hogares totales. & Candidato; proxy de carga potencial de cuidado, no dependencia económica.\\
2HM & $H^{60+}_i/H^{(2HM)}_i$: hogares con personas de 60+ / total de hogares del bloque 2HM. & Candidato; no mide dependencia funcional.\\
6V & $(G_i+T_i+M_i+C_i+D_i+S_i)/V_i$: Guadua, Tapia/bahareque/adobe, Madera burda, Caña/vegetales, Desecho y Sin paredes / viviendas tipo Total. & Candidato; agregado acordado, sin pesos por severidad.\\
8V energía & $E_i^{-}/E_i$: viviendas ocupadas presentes sin energía / total viviendas ocupadas presentes. & Candidato de carencia.\\
8V gas & $G_i^{-}/(G_i^{T}-G_i^{NR})$. & Candidato entre respuestas clasificables; NR se conserva y se sensibiliza.\\
8V internet & $I_i^{-}/(I_i^{T}-I_i^{NR})$. & Candidato entre respuestas clasificables; NR se conserva y se sensibiliza.\\
10V estrato & $(E_{1i}+E_{2i}+N_i)/(\sum_{s=1}^{6}E_{si}+N_i)$. & Candidato acordado; excluye Sin servicio y No sabe/No responde del cociente.\\
5PM actividad & $(U_i+B_i+2C_i)/[2(T_i-N_i)]$: negocio no remunerado, buscó trabajo, incapacidad permanente (peso 2), sobre respuestas válidas ponderadas. & Candidato experimental; no es tasa oficial ni dependencia del hogar.\\
17PM sin media completa & $(Ninguno+Preescolar+Primaria_{c/i}+Secundaria_{c/i}+Media_i)/P_{5+}$. & Descriptiva reservada; fuera del primer PCA.\\
\bottomrule
\end{longtable}
\end{singlespace}

En 5PM el denominador ponderado limita la puntuación a [0,1], bajo categorías exhaustivas y mutuamente excluyentes. El peso 2 para incapacidad es una decisión normativa del equipo, no una estimación del censo. En 17PM las proporciones de cada una de las 15 categorías son $p_{ik}=n_{ik}/P_{5+}$; se controla $|\sum_{k=1}^{15}p_{ik}-1|\leq10^{-10}$. La sensibilidad incorpora Normal incompleta al numerador; Normal completa se excluye. El denominador 5+ incluye Sin información en la fórmula principal, con alternativa que la excluye.

\subsection{Etapa 3. Controles aritméticos, rangos y completitud}

El residuo $e_i$ de la etapa 1 se aplica a los totales publicados y sus categorías: distribución educativa 17PM, actividad 5PM, nueve materiales 6V, energía 8V, gas e internet 8V y categorías 10V. Para la matriz candidata también se verifica cobertura/unicidad, que cada proporción esté en [0,1], y el número de valores NA. La ejecución del 15/10/2026 registró cero violaciones en los controles incluidos y 10.098 celdas completas para nueve candidatas en 1.122 unidades territoriales. Esto confirma coherencia aritmética y cobertura de la salida, no validez teórica del constructo.

En la auditoría del IPM, cada indicador se contrasta contra [0,100], se compara IPM Municipal con IPM, y se reconcilian las hojas ancha y largas por código y variable. Los resúmenes regionales, departamentales y controles se recalculan desde la tabla ancha y se informa su diferencia absoluta. El detalle por prueba está en \path{resultados/verificacion_consolidados/chequeos_consistencia.csv}; las discrepancias no se corrigen automáticamente.

\subsection{Etapa 4. Auditoría descriptiva, correlaciones y extremos}

Para cada variable $x$, se calculan media $\bar{x}=n^{-1}\sum_i x_i$, varianza muestral $s^2=(n-1)^{-1}\sum_i(x_i-\bar{x})^2$, mínimos, máximos y cuantiles tipo 7 de R (1\%, 25\%, 50\%, 75\%, 99\%). La función auxiliar define los momentos estandarizados como
\[
g_1=\frac{n^{-1}\sum_i(x_i-\bar{x})^3}{s^3},\qquad
g_2=\frac{n^{-1}\sum_i(x_i-\bar{x})^4}{s^4}-3,
\]
para asimetría y curtosis en exceso; el script de auditoría también genera esos dos descriptivos con \texttt{psych::describe}, cuya convención de corrección finita puede producir pequeñas diferencias. Se inspecciona porcentaje de ceros $100n^{-1}\sum_i\mathbf{1}(x_i=0)$ y varianza casi nula ($s^2<10^{-8}$). Los umbrales de alerta exploratoria son $|g_1|>2$, al menos 50\% de ceros, o proporciones fuera de [0,1]. No son pruebas de normalidad ni reglas automáticas de exclusión.

La matriz de correlaciones histórica es Pearson: $r_{jk}=\operatorname{cov}(X_j,X_k)/(s_js_k)$, calculada sobre la base total; en la auditoría se recalculó cada par y se comparó con el CSV, con tolerancia $10^{-10}$. Los niveles descriptivos fueron $|r|\geq0,40$ (moderado), $|r|\geq0,70$ (fuerte), y alertas $|r|\geq0,90$ y $|r|\geq0,95$. Estos umbrales no son contrastes de significancia.

La estandarización histórica se define como $z_{ij}=(x_{ij}-\bar{x}_j)/s_j$. Se marcó como extremo si $|z_{ij}|>4$ para al menos un indicador. La matriz previa de siete candidatos produjo 32 unidades con algún valor extremo. Es un indicador para cotejar la fuente, no una razón para corregir o eliminar observaciones. La actualización recalcula 43 códigos con 47 celdas extremas para ocho variables y 1.102 unidades municipales e insulares; el cotejo fuente se documenta en la sección de actualización descriptiva.

La auditoría reproducible de artefactos comprueba que los resultados almacenados coincidan con su recálculo. Para descriptivos se evalúa la máxima diferencia absoluta entre los valores guardados y recalculados de mínimo, percentil 1, mediana, media, percentil 99 y máximo; para correlaciones se usa la máxima diferencia por par. En ambos casos se exige diferencia menor que $10^{-10}$. Para sensibilidad NR se comparan los conteos de municipios marcados y la máxima diferencia absoluta entre resumen y archivo municipal; para extremos se compara el número de municipios con al menos un $|z|>4$. La auditoría histórica obtuvo diferencias de $1,11\times10^{-16}$ en descriptivos, $1,11\times10^{-16}$ en correlaciones y cero en los dos últimos controles: los cuatro controles aprobaron la conciliación de archivos. Esta verificación comprueba reproducción interna, no la validez del indicador.

\subsection{Etapa 5. Sensibilidad de fórmulas y no respuesta}

Para cada municipio se comparan la fórmula principal $a_i$ y la alternativa $b_i$. Se calcula diferencia absoluta $d_i=|a_i-b_i|$, correlación de rangos $\rho_s=\operatorname{cor}(rank(a),rank(b))$ en pares completos, diferencia máxima $\max_i d_i$, y número de unidades con $d_i\geq0,02$. En la auditoría histórica se definió además diferencia relativa $q_i=d_i/|a_i|$ cuando $a_i\ne0$; un municipio era sensible si $d_i\geq0,02$ \textbf{o} $q_i\geq0,10$. El resumen actual de escenarios que acompaña la matriz candidata usa el umbral absoluto de dos puntos porcentuales.

Los escenarios ejecutados son: 5PM con/sin No informa; cota que asigna puntuación máxima a No informa, $(U_i+B_i+2C_i+2N_i)/(2T_i)$; pesos de incapacidad 1:1 y 3:1; gas e internet con denominador válido frente al total publicado; 10V frente al total de personas de BU; y 17PM incluyendo Normal incompleta y excluyendo Sin información del denominador. En 5PM, la alternativa que mantiene numerador observado pero emplea total es $(U_i+B_i+2C_i)/(2T_i)$. En servicios, el escenario alterno divide el conteo sin servicio por el total que incluye Sin información. Los resultados numéricos están en la Tabla \ref{tab:sensibilidad}.

La tasa municipal de NR por ítem es $NR_i=n_{NR,i}/T_i$; la tasa agrupada ponderada es $\sum_i n_{NR,i}/\sum_i T_i$. Se reportan máximo municipal y unidades que exceden los umbrales 5\%, 10\% y 20\%. \enquote{No tiene estrato} es sustantiva, no NR; \enquote{No sabe/No responde} se conserva como categoría no clasificable. Las comparaciones de valores y rangos por municipio no establecen la estabilidad de correlaciones, KMO/MSA, cargas o puntajes.

\subsection{Etapa 6. Adecuación para PCA: fórmulas, resultados y pendientes}

El índice KMO global contrasta la suma de correlaciones ordinarias $r_{ij}$ con las correlaciones parciales $a_{ij}$:
\[
KMO=\frac{\sum_{i<j}r_{ij}^{2}}{\sum_{i<j}r_{ij}^{2}+\sum_{i<j}a_{ij}^{2}}.
\]
El MSA individual para la variable $i$ es $MSA_i=\sum_{j\ne i}r_{ij}^{2}/[\sum_{j\ne i}r_{ij}^{2}+\sum_{j\ne i}a_{ij}^{2}]$. Las correlaciones parciales controlan linealmente las restantes variables. KMO/MSA describen la proporción de correlación simple frente a parcial; son diagnósticos y no una aprobación automática \citep{kaiser1974simplicity}.

La prueba de esfericidad de Bartlett contrasta $H_0:R=I_p$ frente a correlaciones no identidad. El estadístico usado por la implementación es
\[
\chi^2_B=-\left[n-1-\frac{2p+5}{6}\right]\ln|R|,\qquad gl=\frac{p(p-1)}{2},
\]
donde $n$ es el número de unidades, $p$ el número de indicadores y $|R|$ el determinante de la matriz de correlaciones. Un p-valor pequeño evidencia que $R$ difiere de identidad; no demuestra que exista un IVS unidimensional interpretable. Con $n$ grande puede resultar significativo ante asociaciones pequeñas \citep{bartlett1950factor}.

La aproximación inferencial de Bartlett presupone condiciones de muestreo e independencia y un modelo de normalidad multivariada. En agregados censales territoriales, con proporciones acotadas y posible dependencia espacial, no debe interpretarse como una prueba de representatividad ni como aprobación automática del índice. La transformación PCA exploratoria puede calcularse sin exigir normalidad; las garantías inferenciales de un contraste son una cuestión distinta \citep{jolliffe2002principal}.

\begin{table}[htbp]
\centering\small
\caption{Estado de las pruebas de adecuación multivariada}\label{tab:estado_pruebas_pca}
\begin{adjustbox}{max width=\textwidth}
\begin{tabular}{@{}p{0.20\textwidth}p{0.25\textwidth}p{0.40\textwidth}p{0.10\textwidth}@{}}
\toprule
\textbf{Prueba} & \textbf{Matriz} & \textbf{Resultado documentado} & \textbf{Estado}\\
\midrule
KMO global & Histórica, 7 variables & 0,642484; adecuación moderada, no aprobación automática. & Ejecutado\\
Bartlett & Histórica; $n=1122$, $p=7$ & $\chi^2=4664,8728$, $gl=21$; archivo almacena $p=0$ por subdesbordamiento, se informa $p<0,001$. & Ejecutado\\
MSA por variable & Histórica, 7 variables & Estrato 1--2: 0,414311, el MSA más bajo; otros valores en \path{resultados/base_trabajo/kmo_por_indicador.csv}. & Ejecutado\\
KMO/Bartlett & Ocho candidatas; $n=1102$ & KMO 0,672406; $\chi^2=4305{,}6886$, $gl=28$, $p$ nominal $<0,001$; MSA 0,5553--0,7842. & Ejecutado 20/10\\
Sensibilidad de adecuación & Fórmulas y ámbitos alternativos & KMO/MSA ejecutados 20/10; cargas, componentes y puntajes aún diferidos. & Parcial\\
Estabilidad PCA & Fórmulas y universos alternativos & Componentes, cargas, eigenvalores y puntuaciones aún no calculados. & Pendiente\\
\bottomrule
\end{tabular}
\end{adjustbox}
\end{table}

Los resultados KMO/Bartlett disponibles pertenecen a una versión previa de siete candidatos: actividad, energía, gas, recolección de basuras, internet, hogares con menores y estrato 1--2. No corresponden a las nueve candidatas calculadas más recientemente ni a las ocho variables acordadas para el primer PCA. Los 21 grados de libertad son coherentes con $p=7$. El p-valor registrado como cero es una limitación de representación numérica, no probabilidad exactamente nula. El diagnóstico del 20 de octubre ya recalcula estas pruebas para las ocho candidatas y el ámbito municipal e insular; si cambia la especificación final, deben repetirse.

\subsection{Trazabilidad y límites}

Las expresiones y controles anteriores se vinculan con \path{R/02_base_final_indicadores.R}, \path{R/03_auditoria_indicadores.R}, \path{R/05_conformar_indicadores_finales.R}, la configuración \path{R/00_config.R} y sus CSV. Los indicadores y sensibilidades históricos se reproducen desde \path{resultados/indicadores_finales_ivs_20261015/}. Los valores históricos de KMO, Bartlett y MSA pertenecen a la matriz de siete variables; los nuevos resultados de ocho candidatas se documentan por separado en la actualización del 20 de octubre. La actualización descriptiva añade fórmulas, correlaciones y extremos mediante el script 06; no ejecuta PCA. Los chequeos estructurales de los consolidados están desglosados por hoja en \path{resultados/verificacion_consolidados/}. No se ha ejecutado PCA ni se ha validado un IVS final.

